% ============================================================
%  Chapter 7: Conclusion
%  File: chapters/07_conclusion.tex
%  Bibliography: chapters/refs_07_conclusion.bib
% ============================================================

\chapter{Conclusion}
\label{ch:conclusion}

\section{Summary of the Work}

This thesis began with a simple question: \emph{can a neural network
discover its own structure from first principles, without a designer
specifying layers, connectivity, or even the number of nodes?}

We answered this question in the affirmative, through an architecture
we call \modnet{}. The architecture has four defining properties:

\begin{enumerate}
  \item \textbf{Layer-free.} There are no predefined layers, and no
    modular grouping. Every node is a general-purpose threshold unit
    that may connect to any other node, including itself.

  \item \textbf{Outputs as ordinary nodes.} The first $n_{\text{out}}$
    nodes of the network are designated as outputs; there is no separate
    output layer.

  \item \textbf{Grown and pruned.} The network starts from a small random
    topology and grows through addition of nodes, and shrinks through
    removal of inactive or unreachable ones. Growth is bidirectional.

  \item \textbf{Gradient-free.} No gradient information is used at any
    stage. All learning is performed by a compact evolutionary algorithm
    operating on the topology and weights simultaneously.
\end{enumerate}

We evaluated \modnet{} on a suite of \textbf{sixteen} benchmark problems
spanning four broad classes: Boolean logic, digital arithmetic, sorting
networks, and continuous regression. Table~\ref{tab:conclusion-headline}
summarizes the headline results.

\begin{table}[h]
  \centering
  \small
  \caption{Headline results of \modnet{} on the sixteen benchmark
    problems. Boolean tasks are measured in exact accuracy; continuous
    tasks in root-mean-square error (RMSE).}
  \label{tab:conclusion-headline}
  \begin{tabular}{llrrr}
    \toprule
    \textbf{Task} & \textbf{Type} & \textbf{Result} & \textbf{Nodes} & \textbf{Time (s)} \\
    \midrule
    XOR-2               & Boolean    & \textbf{100.0\%} &  8 &  0.8 \\
    XOR-3               & Boolean    & \textbf{100.0\%} &  9 &  2.0 \\
    Parity-4            & Boolean    & \textbf{100.0\%} & 10 &  1.8 \\
    Parity-6            & Boolean    & 95.3\%           &  8 & 40.6 \\
    Majority-3          & Boolean    & \textbf{100.0\%} & 26 &  0.9 \\
    Majority-5          & Boolean    & \textbf{100.0\%} & 26 & 10.2 \\
    Full Adder          & Boolean    & \textbf{100.0\%} & 26 &  3.0 \\
    2-bit Adder         & Boolean    & \textbf{100.0\%} & 26 &  7.3 \\
    Comparator          & Boolean    & \textbf{100.0\%} & 27 &  3.8 \\
    MUX 4-to-1          & Boolean    & \textbf{100.0\%} & 26 &  4.8 \\
    Sort-4              & Boolean    & \textbf{100.0\%} & 26 &  8.7 \\
    \midrule
    $\sin(2\pi x)$      & Regression & RMSE $= 0.0465$  & 18 & 30.0 \\
    $\sin(4\pi x)$      & Regression & RMSE $= 0.0789$  & 13 & 25.1 \\
    $x^2 + y^2$         & Regression & RMSE $= 0.0143$  & 26 & 21.1 \\
    Gaussian            & Regression & RMSE $= 0.0275$  & 27 & 41.6 \\
    Multiplication      & Regression & RMSE $= 0.0254$  & 13 & 28.5 \\
    \bottomrule
  \end{tabular}
\end{table}

\section{What We Learned}

Beyond the raw numbers, this work produced five insights that we
believe are worth carrying forward.

\subsection{Structure Is Discoverable}

The most important lesson is that \emph{structure is discoverable}.
Given only a minimal node model and a fitness function, evolution
found:

- Minimal networks for Boolean tasks (eight nodes for XOR-2).
- Recurrent networks for temporal tasks (up to fourteen self-loops).
- Sparse networks for all tasks ($27\%$--$67\%$ density).
- Networks that solve nontrivial digital circuits exactly.

No human designer specified any of these properties. They emerged from
the interaction between growth, pruning, mutation, and selection.

\subsection{Recurrence Is a Structural Necessity, Not a Design Choice}

We observed that the number of recurrent self-loops scales with the
temporal complexity of the task:

\begin{itemize}
  \item XOR-2 (combinational): one self-loop.
  \item Parity-4 (combinational, deeper): five self-loops.
  \item Gaussian (continuous): fourteen self-loops.
  \item Multiplication (continuous, 2D): ten self-loops.
\end{itemize}

This suggests that recurrence is not a special architectural feature to
be added when needed, but a natural consequence of tasks that require
temporal integration. \modnet{} discovers this necessity without being
told.

\subsection{Sparsity Is Preferred}

Evolved networks use only a fraction of the available connections. The
pruning mechanism removes nodes that are either inactive or saturated,
and the resulting networks have densities of $27\%$ to $67\%$
depending on the task. This is consistent with the sparsity of
biological neural networks \citep{chklovskii2004wiring, cherniak1994component}.

\subsection{Modularity Is Not Necessary}

In earlier work, we studied a modular variant of \modnet{}, where
nodes were grouped into a fixed number of modules. The present work is
fully layer-free: there are no modules at all. A direct comparison
revealed that the layer-free variant is \emph{equal or superior} on
every benchmark, including a decisive win on the multiplication task
($13$ nodes versus $18$). This suggests that modular scaffolds, far
from helping, can restrict the search when the task does not require
modular structure.

\subsection{Digital Circuits Are Within Reach}

To the best of our knowledge, this is the first demonstration of a
fully layer-free evolved network that solves full adder, 2-bit adder,
comparator, 4-to-1 multiplexer, and 4-bit sorting network, all with
exact $100\%$ accuracy. This extends the scope of neuroevolution
beyond classification and control into the domain of digital logic
synthesis.

\section{Limitations}

We do not claim that \modnet{} is a universal solver. Several
limitations must be acknowledged.

\subsection{Scale}

The largest network in our experiments has $27$ nodes. This is small
by modern standards: even a modest feedforward network for a real-world
task typically has thousands of parameters. The primary reason for
this limitation is computational: the fitness evaluation is
$O(E \cdot T)$ per network, and the evolutionary algorithm performs
$O(P \cdot G)$ evaluations, so the total cost grows with the size of
the network.

\subsection{Two Failure Cases}

Two benchmark problems did not reach their success threshold:

\begin{itemize}
  \item \textbf{Parity-6}: stuck at $95.3\%$ accuracy with a network of
    eight nodes. The learning curve shows no sign of approaching exact
    accuracy, suggesting a genuine local optimum.
  \item \textbf{$\sin(4\pi x)$}: reached RMSE $= 0.0789$, close to but
    above the $0.05$ threshold. The network was still improving at the
    end of the run, suggesting that additional generations might
    suffice.
\end{itemize}

These failures are informative: they show that \modnet{} is not a
universal solver, and that certain tasks require either more resources
or a different search strategy.

\subsection{No Theoretical Guarantees}

We have not proven any theoretical properties of \modnet{}. In
particular, we have not shown that the evolutionary algorithm converges
to a global optimum, nor have we formally established Turing
completeness. The empirical results are encouraging, but they do not
constitute a proof.

\subsection{Single Seed}

Our experiments use a fixed random seed. With a different seed, the
algorithm might find different (and possibly better or worse) networks.
A multi-seed analysis would allow us to report means and standard
deviations, and to distinguish robust patterns from coincidences.

\subsection{Limited Task Diversity}

The sixteen tasks we considered all involve low-dimensional inputs
($2$--$8$ bits) and outputs ($1$--$8$ bits). This is far from the
dimensionality of real-world problems. We have not demonstrated that
\modnet{} scales to high-dimensional inputs.

\section{Future Work}

Several directions for future research emerge from this work.

\subsection{Scaling to Larger Networks}

The most obvious direction is to scale \modnet{} to larger networks.
There are several ways to approach this:

\begin{enumerate}
  \item \textbf{Parallelize the fitness evaluation.} Each network in
    the population is independent, so evaluations can be distributed
    across multiple cores or machines. With a cluster of GPUs, a
    population of $10^5$ networks could be evaluated in the time it
    takes to evaluate $10^3$ on a single machine.

  \item \textbf{Sparse encoding.} The current encoding represents each
    weight as a 32-bit float. A sparse encoding that stores only
    active connections could reduce memory and computation by an
    order of magnitude.

  \item \textbf{Advanced search algorithms.} The evolutionary algorithm
    we use is simple. More advanced methods such as CMA-ES,
    differential evolution, or Bayesian optimization might converge
    faster.
\end{enumerate}

\subsection{Hybrid Learning}

A natural extension is to combine evolutionary topology search with
gradient-based weight optimization. In this hybrid approach, the
evolutionary algorithm would discover the topology, and gradient
descent would fine-tune the weights. The challenge is that the
Heaviside activation is not differentiable; one solution is to use a
differentiable surrogate during training, as done in spiking neural
networks \citep{neftci2019surrogate}.

\subsection{Extending to Structured Inputs}

Our current encoding assumes unstructured bit vectors. Extending
\modnet{} to images, sequences, or graphs would require incorporating
structural priors into the architecture, for example through
convolutional or graph-based preprocessing.

\subsection{Theoretical Analysis}

Finally, a more thorough theoretical analysis would be valuable. Key
questions include:

\begin{itemize}
  \item Under what conditions is \modnet{} Turing complete?
  \item What is the convergence rate of the evolutionary algorithm?
  \item Can we characterize the class of tasks that \modnet{} can
    solve efficiently?
\end{itemize}

\section{Broader Implications}

Beyond the technical results, this work has broader implications for
the way we think about neural architecture design.

The dominant paradigm in modern machine learning is to design
architectures by hand, guided by intuition and empirical tuning. This
has been remarkably successful, but it has also led to a proliferation
of architectures that are highly specialized for particular tasks. The
transformer, for example, is optimized for sequences; the
convolutional network is optimized for images; the recurrent network
is optimized for time series. Each architecture comes with its own set
of hyperparameters, its own training procedure, and its own failure
modes.

\modnet{} suggests an alternative: a single, general-purpose
architecture that can be adapted to any task through evolution. The
same node model, the same fitness function, and the same evolutionary
algorithm solved Boolean logic, digital arithmetic, sorting, and
continuous regression with no task-specific modifications.

This is not to say that \modnet{} is better than specialized
architectures on their home turf---it is not, at least not at the
scale we have tested. But it does suggest that architecture design
can be automated, and that the resulting architectures need not be
specialized. This is the essence of meta-learning, and \modnet{} is a
step in that direction.

\section{Final Remarks}

We began this thesis with the observation that biological brains have
no layers, no fixed topology, and no global error signal. They
discover their structure through developmental and evolutionary
processes. \modnet{} is a modest attempt to capture these properties
in a minimal computational model.

Our results suggest that this approach is viable. A single
layer-free architecture, driven by a compact evolutionary process,
solved fourteen of sixteen tasks spanning Boolean logic, digital
arithmetic, sorting networks, and continuous regression, using no more
than $27$ nodes per network. Recurrence emerged spontaneously and
scaled with task complexity. Sparsity emerged without any regularizer.
And a direct comparison with a modular variant showed that removing
structure entirely can \emph{improve} performance.

There is much work still to be done. The architecture has limits---two
of our sixteen tasks were not fully solved---and scaling remains an
open challenge. But the direction is promising. If structure can be
discovered rather than designed, then the field of neural architecture
may be entering a new phase, one in which the designer's role shifts
from \emph{building} networks to \emph{specifying objectives}.

We close with a reflection. The history of machine learning is a
history of \emph{removing} structure from our models. Early models were
heavily hand-crafted. Deep learning removed the feature engineering.
Neural architecture search removed the architecture design. \modnet{}
takes this trend one step further: it removes even the layered
structure itself, and asks the network to discover it.

Perhaps the next step is to remove ourselves entirely from the design
loop. If a network can discover its own structure, perhaps it can
also discover its own objective. That question is beyond the scope of
this thesis, but the tools developed here---self-growing layer-free
networks with complete event logging and reproducible benchmarks---may
serve as a foundation for exploring it.

% ============================================================
%  Bibliography for this chapter
% ============================================================
